\section{Conductor-sensitive shape bounds for all ordinary alternating cores}
\label{sec:ordinary-alternating-conductor}

\paragraph{Status and scope.}
This complete new root derivation has passed the full root and
independent /root/sixfolds proof and code reviews, with separate exact
supplements. The cyclotomic scalar-line footer precision is incorporated. It combines the actual scalar-conductor argument of HD-AST-04--05
with the full-normalizer support estimate of HD-ALT-07. No equality
between the original quadratic field and the inducing character field
is assumed. Its new shape restriction concerns every remaining long
ordinary Specht label, not only the standard label.

\begin{theorem}[HD-ALT-08: arithmetic long-row restriction]
\label{HD-ALT-08}
In the transitive ordinary alternating-core setting of HD-ALT-07, let
$n\ge15$, transpose the label if needed, and write
$\lambda=(n-k,\mu)$ with $1\le k<n/4$.
Suppose that the original imaginary quadratic field $K$ satisfies
one of the following conditions, with the indicated integer $h$:
\[
 h=4\quad\text{if }8\mid|\Delta_K|,
 \qquad h=p-1\quad\text{if an odd prime }p\ge5
                         \text{ divides }|\Delta_K|.
 \tag{AC.1}
\]
If several choices are available, the largest gives the strongest bound.
Every actual rank-nine source in this setting necessarily satisfies
\[
 \beta(n,k):=\frac{k(n-1-k)}{(n-1)(n-2)}
                     \le\frac1{5h}+\frac3{5e},
 \tag{AC.2}
\]
where $e$ is its actual total inducing degree.
For every remaining label after the separate single-standard exclusion,
\[
 \max(\lambda_1,\lambda'_1)>
          \frac{1+\sqrt{1-4/(5h)}}2\,n.
 \tag{AC.3}
\]
Thus the remaining sixfold supplier gate is restricted by the actual
original field, with the same full scalar, outer and tensor-permutation
scope as HD-ALT-07. No source is asserted to exist within the bound.
\end{theorem}

\begin{proof}
First verify the arithmetic interface for these general labels.
Each long Specht module has its ordinary rational $S_n$ model, and
its restriction to $A_n$ is irreducible and nontrivial.
Transitivity identifies the factor labels up to transpose; choosing
the same long-label extension on every factor gives a rational finite
tensor-wreath image $\mathcal R_\lambda$. The normalizer argument of
HD-ALT-03 and HD-ALT-07 writes every actual inducing operator as
$g=\eta T$, with $T\in\mathcal R_\lambda$.
Any different sign choice for an $S_n$ extension is absorbed by $\eta$.

There is no nonidentity scalar in $\mathcal R_\lambda$.
A scalar has trivial action by conjugation on each represented
simple factor, so its tensor-factor permutation is trivial.
For every remaining factor its $S_n$ permutation centralizes $A_n$,
because the represented $A_n$ is faithful and has trivial projective
kernel. That centralizer is trivial, as proved in HD-AST-04.
Hence all factor permutations are identities. The scalar coordinate
is consequently unique and multiplicative, with finite image $\mu_M$.
Rationality gives $F\subseteq\mathbf Q(\zeta_M)$ for the actual
inducing character field $F$. Since $K\subseteq F$, the cited
quadratic conductor criterion gives $|\Delta_K|\mid M$.
This is precisely the criterion and source scope in HD-AST-04,
Milne Chapter V, Remark3.8 and Examples3.10--3.11,
printed pp.~158--159.\cite{HD:MilneCFT2020}

For either choice in (AC.1), the cyclic quotient argument of
HD-AST-04 gives at least two genuine inertia entries whose scalar
orders have totient at least $h$. The characteristic polynomial of
the corresponding actual rational operator $T$ has equal
multiplicities on conjugate primitive roots. Its prescribed scalar
inverse can therefore have eigenspace dimension at most $e/h$.
Each of these two actual inertia entries moves at least
$(1-1/h)e$ dimensions. No trace lower bound for the general Specht
character is needed in this step.

The actual intermediate-cover/period argument supplies at least twelve
genuine inertias and total moving rank $2e+6$. For every other entry,
HD-ALT-07 supplies moving rank at least $\beta(n,k)e$, with its
stronger quarter bound for genuine factor permutations and full
moving rank for a nonidentity scalar. Hence
\[
       2e+6\ge2(1-1/h)e+10\beta(n,k)e.
 \tag{AC.4}
\]
This is exactly (AC.2).

Put
\[
 c=\frac{1-\sqrt{1-4/(5h)}}2,
 \qquad c(1-c)=\frac1{5h},\qquad 0<c<\frac1{16}.
 \tag{AC.5}
\]
For the last inequality, $h\ge4$ gives $c(1-c)\le1/20$,
whereas $(1/16)(15/16)=15/256>1/20$, and the function
$x(1-x)$ is increasing on $[0,1/2]$.
Suppose (AC.3) fails, so $k\ge cn$.

If $k=1$, then $cn\le1$. The single standard factor is already
excluded in the nonsplit source gate, so this remaining case has
at least two tensor factors and $e\ge(n-1)^2$. Therefore
\[
 \beta(n,1)-c(1-c)
 >\frac1{n-1}-\frac1n
 =\frac1{n(n-1)}
 >\frac3{5(n-1)^2}\ge\frac3{5e}.
 \tag{AC.6}
\]
The strict middle comparison is $5(n-1)>3n$, valid for $n\ge15$.
This contradicts (AC.2).

Now let $k\ge2$. The already cited low-degree classification
consequence in HD-ALT-07 gives $e\ge n(n-3)/2$.
If $cn\le1$, monotonicity in $k$ gives
\[
 \begin{split}
 \beta(n,k)-c(1-c)
 &>\frac{2(n-3)}{(n-1)(n-2)}-\frac1n\\
 &>\frac1{2n}>\frac6{5n(n-3)}\ge\frac3{5e}.
 \end{split}
 \tag{AC.7}
\]
Here $2(n-3)/((n-1)(n-2))>3/(2n)$ is equivalent to
$n^2-3n-6>0$, and the next comparison is $5(n-3)>12$.
Both hold throughout the stated range.

Finally suppose $cn\ge1$. The exact difference used in ALT07 is
\[
 \begin{split}
 \beta(n,k)-c(1-c)
 &\ge\frac{c\{(2-3c)n-2+2c\}}{(n-1)(n-2)}\\
 &>\frac{29n-32}{16n^3}
 >\frac6{5n(n-3)}\ge\frac3{5e}.
 \end{split}
 \tag{AC.8}
\]
Indeed $2-3c>29/16$, $-2+2c>-2$, $c\ge1/n$,
and $(n-1)(n-2)<n^2$. The resulting lower bracket is positive.
The final strict rational comparison is equivalent to
$49n^2-595n+480>0$; writing $n=15+x$ gives
$49x^2+875x+2580>0$. Every case contradicts (AC.2), proving (AC.3).
\end{proof}

\paragraph{Explicit original fields and strictness.}
For the class-number-one fields other than the Gaussian and Eisenstein
fields, the values in (AC.1) are respectively
\[
\begin{array}{c|rrrrrrr}
K&\mathbf Q(\sqrt{-2})&\mathbf Q(\sqrt{-7})&
\mathbf Q(\sqrt{-11})&\mathbf Q(\sqrt{-19})&
\mathbf Q(\sqrt{-43})&\mathbf Q(\sqrt{-67})&
\mathbf Q(\sqrt{-163})\\\hline
h&4&6&10&18&42&66&162.
\end{array}
 \tag{AC.9}
\]
The first row threshold is $(1+\sqrt{4/5})n/2$, about
$0.9472135955n$, and the other six are stronger.
For a strict finite-degree witness, take $K=\mathbf Q(\sqrt{-2})$,
$n=100$ and $\lambda=(94,6)$. Its degree is
$\binom{100}{6}-\binom{100}{5}=1116764880$ and its support
fraction is $31/539$. This lies below the old ALT07 threshold
$1/10+3/(5e)$ but above $1/20+3/(5e)$ already at the
single-factor degree, and hence at every tensor degree.
This label is nonstandard and is not a tail-two label; the new
restriction is therefore strict beyond the separate standard and
tail-two tensor statements as well. It is an inequality witness,
not a source construction.

The Gaussian field has no conductor divisor giving $h\ge4$ by this
argument, so its previous shape gate is retained. The Eisenstein
families keep their positive detector bounds. These are restrictions
on actual condition-A suppliers (cote16 scan050; SIGN\_NEXT SN-C03/C04)
in the finite-local-system and weight-one-variation language.
The scalar character defines a line over its cyclotomic field;
restriction of scalars gives a rational representation class.
Its arithmetic conductor does not supply an effective cycle.
Condition I (scan049/cote44 Theorem4), the Cayley--Hamilton projector
of scans048--049, and the stated scheme, EGA/SGA/FGA, $K$-theory,
realization, weight and motivic/$L$-function framework are unchanged.
No new use of Markman, Severn or032 is made.
